\scanpage{38}
% L1–9 | top
\deleted{Soient $A$ et $B$ deux algèbres}

Le produit tensoriel gradué gauche $A\otimes B$ est un foncteur sur la catégorie \inserted{$\unsure{C|\mathcal C}$} des algèbres graduées\inserted{ \unsure{bivariante|bivariantes}}. Ce produit tensoriel (malheureusement !) ne s'interprète pas comme un produit dans la catégorie bivariante : car on n'a pas d'hom de la catégorie $f:A\to A\otimes A$, $f=(f^*,f_*)$, qui jouerait le rôle d'une application diagonale. On trouve néanmoins que c'est un produit \inserted{dans $C$} quand $A,B$ sont toutes deux des Alg. de Poincaré\unsure{\ldots}

% L10–15 | middle
Algèbre bivariante simpliciale : c'est un objet simplicial (\uline{pas} \uline{$1/2$} simplicial !) de la catégorie des Algèbres bivariantes, donc un foncteur \unsure{contravariant|contravar.} de la catégorie des ens finis non vides dans $C$. Il est dit \uline{spécial} ``s'il transforme sommes disjointes en produit tensoriel'', i.e. si

% L16–18 | lower middle
pour $I=I'\amalg I''$, les \deleted{\unsure{hom}}
$A(I')\xrightarrow{\alpha^*}A(I)$, $A(I'')\xrightarrow{\beta^*}A(I)$
[où $\alpha:I'\unsure{\to|-}I$, $\beta:I''-I$) \unsure{induisent|définissent} un \uline{isom}
\[
A(I')\otimes A(I'')\xrightarrow{\sim}A(I'\amalg I'').
\]

% L19–23 | lower
Si \deleted{$A(\Delta^0)$ est une algèbre} on se borne à la sous-catégorie pleine des Alg. biv. simpliciales $A(\unsure{*|\cdot})$ telles que $\unsure{A|\Delta}(\langle\unsure{\omega|0}\rangle)=A_0$ soit une algèbre de Poincaré, on trouve une \uline{équivalence} de la catégorie des Alg. biv. simpliciales et des alg. de Poincaré.
% archive sheet number: 37

